\chapterimage{images/10-rat-za-hiperion.jpg}

\chapter{Rat za Hiperion}

\editorialnote{Ratna situacija, ključne strane sukoba i posledice po planetu i Hegemoniju.}

Dan kada je armada Hegemonije krenula prema Hiperionu počeo je zabavom.

Na više od stotinu pedeset svetova ljudi su slavili odlazak flote. Za većinu stanovnika Mreže rat je još bio udaljeni spektakl: slike brodova, politički govori, zastave i osećaj da učestvuju u nečemu istorijskom.

Na Tau Ceti Centru, sedištu vlasti Hegemonije, slavlje je bilo posebno raskošno.

Hiljade političara, vojnih oficira, bogataša i predstavnika najmoćnijih porodica okupile su se u vrtovima Doma Vlade da posmatraju prolazak armade.

Među njima je bio i Džozef Severn.

Bar se tako predstavljao.

Imao je dvadeset šest godina, kratku smeđu kosu i izgled čoveka koji se još nije sasvim navikao na sopstveno telo.

Zvanično je bio umetnik.

U stvarnosti ni sam nije bio siguran šta je.

Njegovo telo bilo je ljudsko.

Njegova svest bila je rekonstrukcija čoveka koji je umro gotovo hiljadu godina ranije.

Džona Kitsa.

Pesnika čija je nedovršena poema nosila ime \emph{Hiperion}.

Ali Severn nije voleo da sebe naziva Kitsom.

Sećanja koja je nosio pripadala su drugom čoveku.

Ponekad ih je doživljavao kao sopstvena.

Ponekad kao slike iz tuđeg sna.

Zato je prihvatio ime Džozef Severn, po umetniku koji je bio uz pravog Kitsa kada je ovaj umirao u Rimu.

Bio je to identitet dovoljno blizak prošlosti da ga podseća ko je možda bio.

I dovoljno udaljen da ne mora da tvrdi da je zaista taj čovek.

\bigskip

Severn je stajao među gostima dok se nebo iznad Tau Ceti Centra gasilo.

Svetla vrtova bila su prigušena.

Ljudi su prestali da razgovaraju.

Neko je pokazao prema nebu.

``Eno ih.''

Prvo su se pojavili izviđački brodovi.

Njihovi motori ostavljali su kratke svetle tragove.

Zatim su kroz vojni dalekobacač počeli da izlaze bakljobrodovi.

Pa nosači.

Razarači.

Krstarice.

Transporteri sa hiljadama marinaca.

Brodovi za zaštitu perimetra.

I ogromni Skok-brodovi koji će, ako stignu do Hiperiona, otvoriti vojni dalekobacač i povezati udaljeni sistem direktno sa Mrežom.

Nebo se pretvorilo u mrežu crvenih i zlatnih tragova.

Ljudi su počeli da aplaudiraju.

Za nekoliko sekundi hiljade gostiju slavile su rat.

Severn nije.

Posmatrao je lica oko sebe.

Većina tih ljudi nikada nije bila na bojnom polju.

Rat za njih još nije značio krv.

Bio je predstava.

Pitao je jednog čoveka zašto smatra da je sukob neophodan.

Odgovor je bio jednostavan.

Proterani su napali Bresiju.

Sada su došli na Hiperion.

Hegemonija mora da im pokaže da takvo ponašanje neće tolerisati.

Severn je pogledao armadu.

Najveću koncentraciju vojne sile koju je Hegemonija ikada okupila.

Nije osećao sigurnost.

Osećao je nelagodu.

\bigskip

Ubrzo ga je pronašla pripadnica obezbeđenja.

``M. Severn? CEO Gledston može sada da vas primi.''

Meina Gledston čekala ga je u Domu Vlade.

Bila je predsedavajuća Senata i najmoćnija osoba Hegemonije.

Vodila je sistem koji je obuhvatao više od sto milijardi ljudi.

Severn je očekivao političkog diva.

Video je umornu staricu.

Visoku, mršavu, sede kose i prodornih smeđih očiju.

U prostoriji su bili ministri, senatori, generali i jedna projekcija savetnika TehnoCentra po imenu Albedo.

Gledstonova ga je predstavila kao umetnika koji će narednih dana praviti njene portrete.

Nekoliko vojnih oficira očigledno nije razumelo zašto je jedan slikar pozvan na sastanak na kojem se odlučuje o ratu.

Gledstonova se nije obazirala.

``Šta mislite o armadi, M. Severn?''

``Lepa je.''

Jedan general je nezadovoljno frknuo.

``Lepa? Gledate najveću vojnu silu okupljenu u istoriji Hegemonije.''

Gledstonova je nastavila.

``A šta mislite o ratu?''

Severn ju je pogledao.

``Mislim da je glup.''

Prostorija je utihnula.

Podrška ratu u Mreži bila je gotovo potpuna.

Javnost je želela da Hiperion bude odbranjen.

Politička budućnost Gledstonove sada je zavisila od pobede.

``Zašto?'' pitala je.

``Hegemonija nikada nije vodila ovakav rat.''

Generali su odmah počeli da nabrajaju prethodne sukobe.

Pobune.

Građanske ratove.

Bresiju.

Severn je odmahnuo glavom.

To nije bilo isto.

Raniji neprijatelji bili su poznati.

Njihove baze, brodovi, tehnologija i politički ciljevi mogli su se proceniti.

Proterani su vekovima živeli izvan Mreže.

Njihova civilizacija razvijala se odvojeno.

Niko nije pouzdano znao koliko ih ima.

Koliko su tehnološki napredovali.

Šta zapravo žele.

A sada je Hegemonija poslala najveći deo svoje vojne snage tri godine vremenskog duga od najbližeg velikog centra Mreže.

``Osim toga'', rekao je Severn, ``u pitanju je Hiperion.''

Nekolicina prisutnih se nasmešila.

Jedna senatorka ga je pitala da li se plaši Šrajka.

Severn nije odgovorio odmah.

Nije se radilo samo o Šrajku.

Na Hiperionu su postojale Vremenske grobnice.

Predmeti koji su, prema najboljim merenjima, dolazili iz budućnosti.

Postojala su antientropijska polja koja su se upravo urušavala.

Postojao je Šrajk.

A među Grobnicama je sada bilo šestoro hodočasnika i jedna beba.

Ljudi koje je Severn video u snovima.

Ljudi za koje je iz nekog razloga bio vezan.

``Hiperion je nepoznanica'', rekao je.

Čak ni TehnoCentar nije mogao pouzdano da predvidi šta će se tamo dogoditi.

To je bilo ono što ga je plašilo.

\bigskip

Kada su ostali napustili prostoriju, Gledstonova je zadržala Severna.

Sa njom je ostao samo njen najbliži saradnik, Li Hant.

``Zaista mislite da je rat greška?'' pitala je.

``Da.''

Gledstonova se nije uvredila.

Umesto toga postavila mu je pitanje koje ga je iznenadilo.

``Sanjate li još Hiperion?''

Severn je shvatio zbog čega je zapravo pozvan.

Njegova svest bila je povezana sa prvom rekonstruisanom Kitsovom ličnošću.

Ta ličnost je, pre svoje smrti na Lususu, prenela deo sebe u Šrenovu petlju koju je sada nosila Bron Lamija.

Bron se nalazila na Hiperionu.

Zbog nečega što ni Severn ni TehnoCentar nisu razumeli, on je u snovima video ono što se događalo oko nje.

Nije mogao da razgovara sa hodočasnicima.

Nije mogao da utiče na njih.

Ali ih je video.

Ponekad je čak osećao delove njihovih misli i emocija.

Gledstonova je to želela da iskoristi.

Direktna komunikacija sa severom Hiperiona više nije postojala.

Severn je možda bio jedini prozor koji je imala prema Vremenskim grobnicama.

``Poslednji put kada sam ih video'', rekao je, ``svi su bili živi. Osim možda Heta Mastina.''

Ispričao joj je za nestanak templara.

Za Hojtove kruciforme.

Za Silenusovu opsesiju poemom.

Za Kasadovu želju da pronađe Monetu i ubije Šrajka.

Za Sola, koji je nosio Rejčel prema Sfingi dok je dete i dalje svakog dana postajalo mlađe.

Za Bron, koja je izvršavala poslednju želju čoveka čiji deo svesti nosi u sebi.

Na kraju je ispričao i ono što je Konzul priznao.

Da je sarađivao sa Proteranima.

Da je mrzeo Hegemoniju.

I da je upravo on aktivirao uređaj koji je ubrzao otvaranje Vremenskih grobnica.

Gledstonova je prvi put izgledala zaista potreseno.

``Ostali znaju?''

``Da.''

``I nisu ga ubili?''

``Ne.''

``Zašto?''

Severn je razmislio.

``Mislim da više niko od njih ne očekuje da ljudski sud odluči šta će se dogoditi.''

Gledstonova ga je dugo gledala.

``Ostanite u Domu Vlade.''

``Zašto?''

``Zato što želim da znam šta sanjate.''

Severn je shvatio.

Zvanično će biti dvorski umetnik.

Nezvanično će biti njen jedini posmatrač onoga što se događa kod Grobnica.

\bigskip

Dok se na Tau Ceti Centru pripremao rat, hodočasnici su već bili u dolini.

Šestoro odraslih.

Jedna beba.

I Šrajk kojeg nigde nisu mogli da pronađu.

Spustili su se među Grobnice pred zoru.

Čitavog dana išli su od jedne do druge.

Prvo do Sfinge.

Zatim do Grobnice od žada.

Obeliska.

Kristalnog Monolita.

Pećinskih grobnica.

Na kraju prema građevini poznatoj kao Šrajkova Palata.

Svaki put očekivali su da će se nešto dogoditi.

Svaki put dočekivala ih je samo praznina.

Sol je najduže ostao u Sfingi.

U toj građevini Rejčel je pre više od dvadeset godina bila pogođena vremenskom pojavom koja joj je promenila život.

Sada ju je ponovo uneo unutra kao bebu.

Tražio je prostoriju iz njenih zapisa.

Nije je pronašao.

Zidovi su bili prazni.

Stari instrumenti istraživača još su stajali napolju.

Šrajku nije bilo ni traga.

Do večeri strah je počeo da se pretvara u iscrpljenost.

Ulogorili su se kraj Sfinge.

Iznad njih se rat ponovo pojavio kao svetlost na nebu.

Narandžaste i crvene eksplozije prekrivale su zvezde.

``Ko pobeđuje?'' pitala je Bron.

Kasad je pogledao prema svetlosti.

``Proterani ispituju odbranu.''

Silenus je slegnuo ramenima.

``Koga je briga?''

Konzul se nije složio.

Ako Proterani probiju odbranu, Hiperion može biti uništen pre nego što oni uopšte sretnu Šrajka.

Silenus se nasmejao.

``Znači možemo da poginemo pre nego što stignemo do smrti koju smo došli da tražimo. Baš nepristojno.''

Bron mu je rekla da umukne.

Niko više nije imao snage za njegove šale.

\bigskip

Sledećeg jutra Severn je sedeo u Ratnoj sobi Doma Vlade.

Gledstonova je bila za središnjim stolom.

Oko nje su sedeli vojni komandanti, senatori, ministri i savetnici.

Severn je bio pozadi.

Crtao je.

Istovremeno je slušao.

Na velikom taktičkom prikazu lebdeo je sistem Hiperiona.

Planeta.

Meseci.

Položaji flote Hegemonije.

I sve veći broj oznaka koje su predstavljale Proterane.

Prvobitne procene govorile su da će neprijatelj moći da upotrebi nekoliko stotina većih borbenih jedinica.

Novi podaci pokazivali su hiljade.

Jedan roj lansirao je gotovo tri hiljade manjih borbenih letelica samo da bi dovršio jedno opkoljavanje.

Gledstonova je pogledala generale.

``Juče ste govorili o šest ili sedam stotina.''

``Obaveštajni podaci bili su pogrešni.''

``Koliko pogrešni?''

Niko još nije znao.

To je bio prvi ozbiljan znak da Hegemonija možda nije ušla u rat koji razume.

\bigskip

Plan SILE bio je jednostavan samo na mapi.

Glavna flota morala je da zadrži Proterane dovoljno dugo da Skok-brodovi izgrade stabilnu vezu sa vojnom mrežom dalekobacača.

Kada se portal otvori, Hegemonija će moći gotovo trenutno da šalje nove brodove, trupe i municiju.

Bez njega, Hiperion je ostajao udaljen godinama vremenskog duga od pojačanja.

Zato je bitka za nekoliko brodova dalekobacača bila važnija od čitavih eskadrila.

Ako portal bude uspostavljen, Hegemonija može dovesti gotovo neograničene snage.

Ako bude uništen pre otvaranja, flota kod Hiperiona ostaje sama.

Severn je gledao taktičku mapu.

Jedno područje bilo je gotovo prazno.

Sever planete.

Vremenske grobnice nisu čak bile jasno označene kao vojni cilj.

Kao da za generale ne postoje.

A ipak je Severn bio siguran da se upravo tamo nalazi pravi centar rata.

Ne među hiljadama brodova.

Nego među šestoro ljudi i jednom bebom.

\bigskip

Dok su oficiri govorili o orbitama, municiji i procenama gubitaka, na samom Hiperionu rat je već imao drugačije lice.

Milioni ljudi pobegli su prema Kitsu.

Polja oko prestonice pretvorena su u izbegličke logore.

Šume su posečene za ogrev i skloništa.

Potoci su bili zagađeni.

Vojne kolone kretale su se putevima između šatora.

Deca su trčala bosa po blatu.

Ljudi su stajali u redovima za vodu i hranu.

Neki su bežali od Proteranih.

Drugi od Šrajka.

Većina od oba.

Guverner-General Teo Lejn pokušavao je da održi grad funkcionalnim.

Njegovi prioriteti bili su jasni.

Evakuacija.

Odbrana od Proteranih.

I sprečavanje panike zbog Šrajka.

Problem je bio što nijedan od ta tri zadatka nije imao dobro rešenje.

Evakuacija miliona ljudi bila je gotovo nemoguća bez stabilnog dalekobacača.

SILA nije želela da oslabi vojnu odbranu korišćenjem transporta za civile.

A niko nije znao kako da zaustavi Šrajka.

Čitava jedna jedinica teritorijalne odbrane već je otišla prema severu.

Nije se vratila.

\bigskip

Na Tau Ceti Centru javnost je i dalje čekala govor Meine Gledston.

Do tada je rat formalno još mogao da se nazove vojnom operacijom.

To će uskoro prestati.

Severn je izašao iz Ratne sobe pre završetka brifinga.

U hodnicima, vrtovima i na terminalima ljudi su već slušali vesti.

Na svim svetovima Mreže život se nastavljao.

Prodavnice su bile otvorene.

Brodovi su plovili rekom Tetis.

Ljudi su išli na posao.

Deca u škole.

Ali gotovo svako je istovremeno bio povezan sa SveStvari i čekao da čuje šta će reći žena koja govori u ime Hegemonije.

Gledstonova se pojavila na milijardama ekrana.

Njen glas prošao je kroz datasferu i fetlinijske veze.

Objavila je ono što su svi već znali.

Hegemonija je u ratu.

Nije obećala brz završetak.

Nije pokušala da ubedi ljude da neće biti žrtava.

Tražila je jedinstvo.

Rekla je da će Hegemonija voditi rat svom snagom koju poseduje.

I da je cilj samo jedan.

Pobeda.

Jer bez pobede možda neće biti opstanka.

\bigskip

Ljudi su slušali na stotinama svetova.

Na nekima su aplaudirali.

Na drugima su samo ćutali.

Većina još nije mogla da zamisli kako bi rat kod jednog udaljenog sveta mogao da promeni njihove živote.

Severn je mogao.

Ne zato što je bio vojnik.

Nije bio.

Ne zato što je znao planove TehnoCentra.

Nije ih znao.

Nego zato što je svake noći sanjao Hiperion.

Video je hodočasnike kako sede kraj vatre.

Video je Sola kako drži svoju pet dana staru kćer.

Hojta koji jedva podnosi bol dve kruciforme.

Kasada koji noću stražari sa puškom u rukama.

Silenusa koji pokušava da se podsmeva strahu.

Bron koja oseća da je previdela nešto važno.

Konzula koji je svojim postupcima pomogao da se sve pokrene.

A oko njih su stajale građevine koje su dolazile iz budućnosti.

Na nebu iznad njih ljudi su ratovali zbog Hiperiona.

Ali niko još nije znao za šta se zapravo bore.

\bigskip

Te noći Severn je ponovo sanjao.

Vetar se pojačavao u Dolini Vremenskih grobnica.

Pesak je prekrivao logor.

Sfinga je svetlucala od statičkog elektriciteta.

Bron se probudila.

Hojta nije bilo u šatoru.

Kasad ga je video kako ide prema Grobnicama.

Negde u oluji nestajala je prilika u crnom plaštu.

Bron je uzela oružje.

``Idem za njim.''

Kasad je ostao da čuva ostale.

Bron je krenula sama.

Pesak je brisao tragove gotovo istog trenutka kada bi nastali.

Prošla je pored Sfinge.

U daljini je svetlela Grobnica od žada.

Na trenutak je videla tamnu figuru kako nestaje u njenom ulazu.

Požurila je za njom.

A na drugom kraju Mreže Džozef Severn spavao je u sobi Doma Vlade i video sve.

Nije mogao da je upozori.

Nije mogao da joj pomogne.

Mogao je samo da sanja.

Rat je počeo.

Ali ono najvažnije na Hiperionu tek je počinjalo.